%%%%%%%%%%%%%%%%%%%%%%%%%%%%%%%%%%%%%%%%%%%%%%%%%%%%%%%%%%%%%%%%%%%%%%%%%%%%%%%%%
%	LITERATURE REVIEW
%   Mandated by Project Task 4.2 ("literature review on Genetic Algorithms and
%   constraint handling techniques").
%%%%%%%%%%%%%%%%%%%%%%%%%%%%%%%%%%%%%%%%%%%%%%%%%%%%%%%%%%%%%%%%%%%%%%%%%%%%%%%%%
\lhead{Literature Review}
\chapter{Literature Review}
\label{SecLiteratureReview}

The Traveling Salesman Problem (TSP) is a standard benchmark for
combinatorial optimisation. The objective is simple but the search space
grows factorially, which makes exact solvers impractical for repeated
experimentation~\cite{applegate2006traveling,reinelt1991tsplib}. Genetic
Algorithms maintain a population of candidate tours rather than a single
incumbent, making them a natural fit for exploring permutation
spaces~\cite{holland1975adaptation,goldberg1989genetic}. Their effectiveness
on the TSP, however, depends on how crossover and mutation handle the
permutation constraint and what happens when infeasible offspring appear.

\section{Genetic Algorithms for the TSP}
\label{Lit:GAforTSP}

The standard representation encodes a tour as a permutation of city
indices. A chromosome is feasible only when every city appears exactly
once. Larra{\~n}aga et al. survey several encodings---path, adjacency,
ordinal, matrix---and the operators designed for
them~\cite{larranaga1999genetic}. Later reviews identify crossover design,
mutation design, local search hybridisation, and diversity preservation as
the recurring themes in GA-TSP
research~\cite{scholz2019historicalreview,cicirello2023permutationoperators}.

Standard string crossover can destroy the city set, so
permutation-specific operators were developed early on. Partially Mapped
Crossover (PMX) preserves positional mapping, Order Crossover (OX) preserves
relative order, and Cycle Crossover (CX) preserves cycles of position-wise
inheritance~\cite{goldberg1985alleles,davis1985applying,oliver1987study}.
Cicirello classifies permutation operators by what they transmit: absolute
positions, undirected edges, directed edges, precedences, or cyclic
precedences~\cite{cicirello2023permutationoperators}. Since TSP tour length
depends on adjacent city pairs, positional diversity---measured by Hamming
distance---captures how the population spreads across the city-position
space.

State-of-the-art GA variants go well beyond a plain generational loop. Liu
uses edge swapping and local search on large
benchmarks~\cite{liu2014powerfulga}. Eremeev and Kovalenko show that
recombination design, neighbourhood mutation, restarts, and diversity
management all matter for asymmetric
instances~\cite{eremeev2017optimalrecombination}. These are context rather
than direct baselines; the work here stays focused on GA fundamentals and
repair on kroA100.

\section{Constraint Handling in Evolutionary Algorithms}
\label{Lit:ConstraintHandling}

Coello Coello surveys constraint-handling techniques for evolutionary
algorithms: penalty functions, specialised representations, repair methods,
and feasibility-based selection~\cite{coello2002constraint}. Michalewicz and
Schoenauer argue that constraints are part of the problem structure, not an
implementation afterthought~\cite{michalewicz1996evolutionary}. For the TSP,
the ``visit every city exactly once'' condition is not a side constraint---it
defines what makes a chromosome a valid tour.

Penalty functions leave infeasible solutions in the population but add a
cost proportional to constraint violation. The tuning problem is
well-known: too weak and the search rewards invalid short tours, too strong
and useful partial structure is discarded. Chehouri et al. propose a
violation-based alternative to avoid
penalty-parameter tuning altogether~\cite{chehouri2016violationfactor}.

A different approach prevents infeasibility by construction. PMX, OX, and CX
are examples: every offspring is automatically a valid tour. The trade-off is
that restricting operators to feasibility-preserving ones makes it harder to
study what happens when infeasible offspring do appear.

\section{Repair Mechanisms in Permutation-based EAs}
\label{Lit:Repair}

Instead of discarding or penalising infeasible chromosomes, repair modifies
them into feasible candidates before evaluation. Padhye et al. note an
important caveat: repair restores feasibility but can also bias the search
if repaired solutions cluster in particular regions of the feasible
space~\cite{padhye2015feasibilitypreserving}. Salcedo-Sanz treats repair as
a broad class of constraint-handling techniques rather than simple
preprocessing~\cite{salcedosanz2009repair}.

For the TSP, naive crossover can produce duplicated and missing cities. A
repair mechanism scans the chromosome, removes duplicates, and inserts
missing cities. How those cities are inserted matters. Random insertion
preserves diversity but disturbs edges. Positional insertion retains more
chromosome structure but reinforces parent bias. Nearest-neighbour insertion
improves local quality but narrows convergence.

Feasibility and quality are different properties. A repaired chromosome is
valid but not necessarily a good tour---the repair may alter the building
blocks inherited from parents. Uray et al. reinforce this point by showing
that tour symmetries and edge structure affect whether fit parents produce
fit offspring~\cite{uray2023csrx}. Repair should therefore be evaluated by
its effect on solution quality, diversity, and convergence, not just
feasibility.

\section{Comparison Studies of TSP Operators}
\label{Lit:CrossoverStudies}

Abdoun and Abouchabaka compare several crossover operators for the TSP,
including OX, PMX, UPMX, CX, and NWOX, reporting strong performance for OX
on their benchmark~\cite{abdoun2012adaptivecrossover}. A related study by
Abdoun, Abouchabaka, and Tajani covers mutation operators and emphasises
that representation, selection, crossover, mutation, and operator
probabilities jointly define a GA
variant~\cite{abdoun2012mutationoperators}.

Thanh et al. test combined crossover mechanisms on TSPLIB instances from 51
to 575 cities, including kroA100 and
lin105~\cite{thanh2020combinationcrossover}. Uray et al. design crossover
operators around symmetric TSP invariances such as circular shifts and tour
reversal~\cite{uray2023csrx}. Both studies show that crossover choice
affects mean and best tour quality on standard benchmarks.

Taken together, the literature points to three things that matter for
GA-based TSP solvers: permutation-aware operators, careful handling of
infeasible offspring, and critical evaluation of repair---since restoring
feasibility is not the same as preserving useful structure.
